% TikZ 图：斜入磁场的立体螺旋运动
% Day 22：洛伦兹力
\begin{tikzpicture}[
  scale=0.9,
  >=Latex,
  line cap=round,
  line join=round,
  x={(1cm,0cm)},
  y={(0cm,1cm)},
  z={(0.34cm,0.20cm)}
]
  \path[use as bounding box] (-3.65,-1.85) rectangle (3.85,1.80);

  % 左侧：入射速度分解
  \begin{scope}[shift={(-3.15,0)}]
    \coordinate (P) at (0,0);
    \fill[red] (P) circle (2pt);
    \draw[->,orange!90!black,very thick] (P) -- (1.05,0.82);
    \node[orange!90!black,above] at (0.72,0.60) {$v$};
    \draw[->,orange!90!black,thick] (P) -- (1.05,0);
    \node[orange!90!black,below] at (0.55,0) {$v_\parallel$};
    \draw[->,orange!90!black,thick] (1.05,0) -- (1.05,0.82);
    \node[orange!90!black,right] at (1.05,0.42) {$v_\perp$};
  \end{scope}

  % 右侧：沿磁场方向推进的立体螺旋线
  \draw[->,blue!70!black,very thick]
    (-1.45,0,0) -- (3.15,0,0)
    node[right] {$B$};
  \node[blue!70!black,above] at (1.05,0.72,0)
    {螺旋轴线 $\parallel B$};

  % 起、末截面帮助读者看出轨迹绕轴旋转
  \draw[gray!55,thin,domain=0:360,samples=80,variable=\t]
    plot ({-1.25},{0.52*cos(\t)},{0.52*sin(\t)});
  \draw[gray!55,thin,domain=0:360,samples=80,variable=\t]
    plot ({2.49},{0.52*cos(\t)},{0.52*sin(\t)});

  \draw[red,very thick,->,domain=0:720,samples=220,variable=\t]
    plot ({-1.25+0.0052*\t},{0.52*cos(\t)},{0.52*sin(\t)});

  % 半径与螺距
  \draw[->,orange!90!black,thick]
    (0.62,0,0) -- (0.62,0.52,0)
    node[midway,left] {$r$};
  \draw[<->,purple,thick]
    (-0.32,-0.85,0) -- (1.55,-0.85,0)
    node[midway,below] {$h=v_\parallel T$};
  \draw[gray,dashed] (-0.32,-0.72,0) -- (-0.32,0.52,0);
  \draw[gray,dashed] (1.55,-0.72,0) -- (1.55,0.52,0);

  \node[font=\small] at (0,-1.55)
    {沿磁场方向匀速前进 $+$ 垂直磁场方向圆周运动 $=$ 螺旋运动};
\end{tikzpicture}
